\section{Background and Related Work}

\subsection{The Ge'ez Script}
The Ge'ez script, known natively as \textit{fidäl}, stands as the only indigenous writing system of the African continent to have evolved continuously into a modern, productive form. With its earliest attested inscriptions traced to the 4\textsuperscript{th} century CE, exemplified by the Hawulti obelisk in Matara, Eritrea, the script emerged within the epigraphic traditions of the Aksumite civilization of the northern Horn of Africa and represents a sophisticated adaptation of earlier South Semitic conventions~\cite{phillipson2012foundations,getatchew1996ethiopic}. Unlike the alphabetic systems of Greek or Latin, Ge'ez is fundamentally a syllabic abugida: each character encodes a consonant-vowel unit~\cite{getatchew1996ethiopic}. This structural innovation arose through deliberate transformation of the South Arabian Sabean script, which Ge'ez adopted and modified by introducing a vowel modification system and shifting the writing direction from right-to-left to left-to-right~\cite{getatchew1996ethiopic,gebremedhin2020linguistic}. This morphological complexity and script-level unity across multiple languages (Amharic, Tigrinya, Tigre, and Classical Ge'ez) create a unique computational challenge.

Amharic adopted the base Ge'ez characters and added characters for additional phonemes, including Cushitic sounds encoded via modified Ge'ez forms. Tigrinya retained the full Ge'ez character set and added eight additional \textit{fidels}~\cite{gebremedhin2020linguistic}. These script properties (large character inventories, morphosyllabic structure, and limited presence in multilingual corpora) directly motivate the need for language-specific tokenization.

\subsection{Vocabulary Expansion in Multilingual Models}
Vocabulary expansion methods aim to augment the tokenizer and embedding matrix of a pretrained multilingual model to better serve target languages underrepresented during pretraining. \citet{wang2019improving} identified fixed vocabulary size as a structural bottleneck in multilingual models and proposed vocabulary adaptation strategies. More recently, \citet{kim2024efficient} demonstrated vocabulary expansion for Korean (EEVE-Korean) via parameter freezing and subword embedding initialization, and reported strong downstream performance. \citet{yamaguchi2024how} examined the minimal data requirements for vocabulary expansion, and \citet{hong2024accelerating} proposed language-specific model heads with customized vocabulary sets combined with verification-based fine-tuning.

\citet{remy2024trans} introduced Trans-Tokenization, a framework for repurposing high-resource models for low-resource settings via adaptive tokenization. \citet{chelombitko2024qtok} presented QTok, a comprehensive framework for evaluating multilingual tokenization quality across languages.

This work differs from these efforts in two key respects. First, we focus on Ge'ez-script languages, which have received little systematic attention in the vocabulary expansion literature despite their significant representational challenges. Second, we combine vocabulary expansion with a two-stage pretraining and fine-tuning pipeline and isolate the contribution of each stage in an ablation.

\subsection{Tokenization for Morphologically Rich Languages}
Subword tokenizers (WordPiece, BPE, SentencePiece) reduce vocabulary size while maintaining generalization across common forms. However, in low-resource settings, particularly for morphologically rich languages, these tokenizers suffer from over-segmentation, high OOV rates, and data inefficiency~\cite{limisiewicz2023tokenization,beinborn2023analyzing}. The problem is amplified for languages with complex derivational or inflectional morphology~\cite{chai2024tokenization}.

Vocabulary-free approaches offer an alternative: ByT5~\cite{xue2022byt5} and CANINE~\cite{clark2022canine} operate at the byte or character level and demonstrate robustness across many languages. However, these models incur higher computational costs due to longer sequences and slower inference~\cite{xue2021mt5,liang2023xlmv}. \citet{pagnoni2024byte} proposed the Byte Latent Transformer (BLT), which processes raw bytes without tokenization, representing a promising but computationally demanding alternative.

\subsection{Multilingual Models for Languages in Africa}
Most multilingual modeling research has concentrated on high-resource languages. Some efforts have expanded coverage: \citet{imanigooghari2023glot500} proposed Glot500-m, which covers 511 predominantly low-resource languages, significantly expanding the set of supported languages. \citet{alabi2022adapting} fine-tuned XLM-R for languages in Africa. AfroXLMR~\cite{alabi2022adapting} and AfriBERTa~\cite{ogueji2021small} provide specialized encoders for some languages in Africa.

Despite this progress, Ge'ez-script languages such as Amharic and Tigrinya remain underserved in terms of tokenization optimization and vocabulary coverage. This work directly targets this gap by constructing language-specific tokenizers, augmenting XLM-R's vocabulary with Ge'ez-script (Amharic and Tigrinya) subword tokens, and evaluating the resulting model on Amharic and Tigrinya downstream tasks.
